\documentclass[10pt,a4paper]{article}

% ==============================================================================
% TEMPLATE EDITORIAL — PERIÓDICO EM 2 COLUNAS (DOSSIÊ ANUAL E CARTA DE CONJUNTURA)
% JORNAL ECONÔMICO DO TOCANTINS
%   - Dossiê Anual (DA): Verde-Floresta (#0B5D32), Marfim (#FAF6EE) e Carmim (#A61C1C)
%   - Carta de Conjuntura (CC): Azul Cobalto (#0A2E73), Esmeralda (#1B7A43),
%     Lima Editorial (#8CE000 no selo escuro) e Azul-Gelo (#F3F6FB)
% ==============================================================================
\usepackage[utf8]{inputenc}
\usepackage[T1]{fontenc}
\usepackage[brazil]{babel}
\usepackage{helvet}
\renewcommand{\familydefault}{\sfdefault}

\usepackage[top=1.35cm, bottom=1.45cm, left=1.35cm, right=1.35cm, headheight=14pt, headsep=8pt, footskip=18pt]{geometry}
\usepackage{xcolor}
\usepackage{graphicx}
\usepackage{booktabs}
\usepackage{tabularx}
\usepackage{longtable}
\usepackage{array}
\usepackage{colortbl}
\usepackage{tikz}
\usepackage{fancyhdr}
\usepackage{lastpage}
\usepackage{enumitem}
\usepackage{microtype}
\usepackage{setspace}
\usepackage{multicol}
\usepackage{amssymb}
\usepackage{hyperref}

% Paleta Editorial Dinâmica (Proporção 75% Neutro / 20% Primária / 5% Acento)
\definecolor{corprimaria}{HTML}{0F2537}
\definecolor{corsecundaria}{HTML}{334155}
\definecolor{coracento}{HTML}{9A6C11}
\definecolor{coracentoheader}{HTML}{C8942B}
\definecolor{coracentotexto}{HTML}{0F2537}
\definecolor{corfundo}{HTML}{FAF8F5}
\definecolor{cinzatexto}{HTML}{1E293B}
\definecolor{cinzasub}{HTML}{475569}
\definecolor{cinzaborda}{HTML}{CBD5E1}
\definecolor{verdealta}{HTML}{15803D}
\definecolor{vermelhobaixa}{HTML}{B91C1C}

\hypersetup{
  colorlinks=true,
  linkcolor=corprimaria,
  urlcolor=corsecundaria,
  pdftitle={Dossiê Anual — Demografia Empresarial, Simples Nacional e Quadro Societário (QSA)},
  pdfauthor={Jornal Economico do Tocantins}
}

\setlength{\parindent}{0pt}
\setlength{\parskip}{3.5pt}
\setlength{\columnsep}{0.60cm}
\setlength{\columnseprule}{0.4pt}
\def\columnseprulecolor{\color{cinzaborda}}
\renewcommand{\arraystretch}{1.14}

% Cabeçalho e Rodapé de Periódico Econômico (Sem sobreposição central)
\pagestyle{fancy}
\fancyhf{}
\renewcommand{\headrulewidth}{0pt}
\renewcommand{\footrulewidth}{0.5pt}
\renewcommand{\footrule}{\hbox to\headwidth{\color{coracento}\leaders\hrule height \footrulewidth\hfill}}

\fancyhead[L]{\ifnum\value{page}>1 \scriptsize\bfseries\color{corprimaria}JORNAL ECONÔMICO DO TOCANTINS \enspace{\color{coracento}$\vert$}\enspace {\mdseries\color{cinzasub}DOSSIÊ ANUAL / SEMESTRAL (DA)}\fi}
\fancyhead[R]{\ifnum\value{page}>1 \scriptsize\color{cinzasub}Ref.: \textbf{Acumulado em 12 Meses (Novembro/2025 a Outubro/2026)} $\cdot$ \#06/2026\fi}

\fancyfoot[L]{\scriptsize\color{cinzasub}\textbf{\color{corprimaria}Jornal Econômico do Tocantins} $\cdot$ Lorenzo Costa | Graduado em Economia — Jornal Econômico do Tocantins $\cdot$ Publicado em 14/11/2026}
\fancyfoot[R]{\scriptsize\bfseries\color{corprimaria}Pág. \thepage\ de \pageref*{LastPage}}

% Banner Editorial Limpo de Abertura de Seção Principal (Caderno Periódico)
\newcommand{\bannersecaoperiodico}[3]{%
  \par\vspace{0.16cm}%
  \noindent{\color{corprimaria}\hrule height 1.4pt}%
  \par\vspace{4pt}%
  \noindent{\scriptsize\bfseries\color{coracento}SEÇÃO #1 \enspace$\cdot$\enspace}{\normalsize\bfseries\color{corprimaria}#2}%
  \par\vspace{2pt}%
  {\scriptsize\itshape\color{cinzasub}#3\par}%
  \vspace{3pt}%
  {\color{cinzaborda}\hrule height 0.5pt}%
  \par\vspace{0.10cm}%
}

% Subtítulo de Coluna dentro do Layout de 2 Colunas
\newcommand{\secaoperiodico}[1]{%
  \par\vspace{0.10cm}%
  \noindent{\small\bfseries\color{corprimaria}#1}%
  \par\vspace{0.03cm}%
  {\color{coracento}\hrule height 0.6pt}%
  \vspace{0.06cm}%
}

\begin{document}
\color{cinzatexto}

% ==============================================================================
% MASTHEAD EDITORIAL CLARO DE PERIÓDICO (DOSSIÊ ANUAL OU CARTA DE CONJUNTURA)
% ==============================================================================
\noindent{\color{corprimaria}\hrule height 2.2pt}%
\par\vspace{2pt}%
{\color{coracento}\hrule height 0.8pt}%
\par\vspace{6pt}%
\noindent\hbox to \linewidth{%
  {\scriptsize\bfseries\color{coracento}\MakeUppercase{DOSSIÊ ANUAL / SEMESTRAL (DA)} \enspace$\cdot$\enspace \#06/2026}%
  \hfill
  {\scriptsize\color{cinzasub}\textbf{\color{corprimaria}JORNAL ECONÔMICO DO TOCANTINS} \enspace$\vert$\enspace Publicado em 14/11/2026}%
}%
\par\vspace{5pt}%
{\fontsize{14.5}{17.4}\selectfont\bfseries\color{corprimaria}Dossiê Anual — Demografia Empresarial, Simples Nacional e Quadro Societário (QSA)\par}%
\vspace{3pt}%
{\footnotesize\itshape\color{cinzasub}Balanço de 12 meses das aberturas, baixas e estoque de CNPJs ativos por porte (MEI, ME e EPP), adesão ao Simples Nacional, perfil societário (QSA) e eficiência da REDESIM no Tocantins \enspace$\vert$\enspace \textbf{\color{corprimaria}Ref.: Acumulado em 12 Meses (Novembro/2025 a Outubro/2026)}\par}%
\par\vspace{6pt}%
{\color{cinzaborda}\hrule height 0.6pt}%

\vspace{0.12cm}

% ==============================================================================
% CORPO COMPLETO DO PERIÓDICO (DOSSIÊ ANUAL OU CARTA DE CONJUNTURA)
% ==============================================================================
\noindent\makebox[\textwidth][s]{%
\begin{tikzpicture}[baseline=(box.base)]
  \node[fill=white, draw=cinzaborda, line width=0.35pt,
        inner xsep=6pt, inner ysep=4pt, text width=0.206\textwidth, minimum height=1.14cm, align=left] (box) {%
    {\fontsize{6.0}{7.2}\selectfont\bfseries\color{cinzasub}CNPJS ATIVOS (ESTOQUE)}\par\vspace{1.2pt}%
    {\fontsize{9.2}{10.8}\selectfont\bfseries\color{corprimaria}167.565}\par\vspace{1pt}%
    {\fontsize{5.8}{6.8}\selectfont\color{cinzasub}Saldo 12m: + 15.608}%
  };
  \fill[coracento] (box.north west) rectangle ([yshift=-1.6pt]box.north east);
\end{tikzpicture}\hfill%
\begin{tikzpicture}[baseline=(box.base)]
  \node[fill=white, draw=cinzaborda, line width=0.35pt,
        inner xsep=6pt, inner ysep=4pt, text width=0.206\textwidth, minimum height=1.14cm, align=left] (box) {%
    {\fontsize{6.0}{7.2}\selectfont\bfseries\color{cinzasub}MEI / MICROEMPRESAS}\par\vspace{1.2pt}%
    {\fontsize{9.2}{10.8}\selectfont\bfseries\color{corprimaria}52,9\% / 34,0\%}\par\vspace{1pt}%
    {\fontsize{5.8}{6.8}\selectfont\color{cinzasub}Pequenos Negócios: 86,9\%}%
  };
  \fill[coracento] (box.north west) rectangle ([yshift=-1.6pt]box.north east);
\end{tikzpicture}\hfill%
\begin{tikzpicture}[baseline=(box.base)]
  \node[fill=white, draw=cinzaborda, line width=0.35pt,
        inner xsep=6pt, inner ysep=4pt, text width=0.206\textwidth, minimum height=1.14cm, align=left] (box) {%
    {\fontsize{6.0}{7.2}\selectfont\bfseries\color{cinzasub}OPTANTES SIMPLES (SINAC)}\par\vspace{1.2pt}%
    {\fontsize{9.2}{10.8}\selectfont\bfseries\color{corprimaria}164.996}\par\vspace{1pt}%
    {\fontsize{5.8}{6.8}\selectfont\color{cinzasub}Sócios QSA: 106.471}%
  };
  \fill[coracento] (box.north west) rectangle ([yshift=-1.6pt]box.north east);
\end{tikzpicture}\hfill%
\begin{tikzpicture}[baseline=(box.base)]
  \node[fill=white, draw=cinzaborda, line width=0.35pt,
        inner xsep=6pt, inner ysep=4pt, text width=0.206\textwidth, minimum height=1.14cm, align=left] (box) {%
    {\fontsize{6.0}{7.2}\selectfont\bfseries\color{cinzasub}TEMPO DE ABERTURA (2026)}\par\vspace{1.2pt}%
    {\fontsize{9.2}{10.8}\selectfont\bfseries\color{corprimaria}18,5 horas}\par\vspace{1pt}%
    {\fontsize{5.8}{6.8}\selectfont\color{cinzasub}Viab.: 15,7h | Reg.: 2,7h}%
  };
  \fill[coracento] (box.north west) rectangle ([yshift=-1.6pt]box.north east);
\end{tikzpicture}%
}

\vspace{0.03cm}
\par\vspace{0.05cm}
\noindent{\color{cinzaborda}\hrule height 0.5pt}\par\vspace{4pt}
{\fontsize{7.8}{9.4}\selectfont\bfseries\color{corprimaria}\MakeUppercase{Síntese dos Indicadores — Balanço em 12 Meses}}\par\vspace{2pt}
{\scriptsize\color{cinzatexto}
\begin{itemize}[leftmargin=10pt, itemsep=1.2pt, topsep=1pt, parsep=0pt]
   \item \textbf{\color{corprimaria}Estoque de CNPJs Ativos (Out/2026):} Estoque de CNPJs Ativos (Out/2026) registrou aceleração expressiva de +9,65\% (mudança drástica acima do padrão recente). (patamar atual: 167.565)
   \item \textbf{\color{corprimaria}Novas Inscrições de CNPJ no Mês:} Novas Inscrições de CNPJ no Mês registrou aceleração expressiva de +5,46\% (mudança drástica acima do padrão recente). (patamar atual: 2.760)
   \item \textbf{\color{corprimaria}Participação de MEIs + Microempresas (ME):} Participação de MEIs + Microempresas (ME) manteve-se praticamente inalterado (+0,05\%), indicando estabilidade marginal. (patamar atual: 86,9\%)
   \item \textbf{\color{corprimaria}Tempo Médio de Abertura de Empresas (TO vs. Brasil):} Tempo Médio de Abertura de Empresas (TO vs. Brasil) sofreu retração acentuada de -29,39\% (queda drástica que requer atenção conjuntural). (patamar atual: 18,5h (BR: 26,2h))
\end{itemize}
}
\noindent{\color{cinzaborda}\hrule height 0.5pt}\par\vspace{0.05cm}

\vspace{0.03cm}
\begin{multicols}{2}
\secaoperiodico{1. Análise Estrutural Anual / Acumulado em 12 Meses}
{\small A avaliação do ciclo de 12 meses da dinâmica empresarial no Tocantins revela uma expansão líquida robusta de 15.608 estabelecimentos ativos, equivalente a um avanço superior a 10\% no estoque empresarial do Estado em apenas um ano. Em todos os doze meses da janela analisada, o volume de novas inscrições superou amplamente o número de baixas cadastrais, confirmando um ambiente favorável ao nascimento e sobrevivência de negócios.

\medskip
Os pequenos negócios constituem a espinha dorsal dessa expansão: somados, os Microempreendedores Individuais (52,9\%) e as Microempresas (34,0\%) representam 86,9\% das empresas ativas no Estado, às quais se somam 7,1\% de Empresas de Pequeno Porte (EPP). No Quadro de Sócios e Administradores (QSA), o predomínio da figura do Sócio-Administrador (86,7\% das qualificações) e a concentração etária entre 31 e 50 anos atestam a gestão direta dos proprietários sobre os empreendimentos locais.}
\end{multicols}
\begin{center}\includegraphics[width=0.84\textwidth]{figuras/2026_11_dossie_empresas_emp_ativos.pdf}\end{center}
\par\vspace{0.08cm}
\noindent\begin{minipage}{\linewidth}
  \noindent{\fontsize{7.8}{9.5}\selectfont\bfseries\color{corprimaria}\rule{0.10cm}{0.20cm}\hspace{4pt}Distribuição das Empresas Ativas por Porte no Tocantins (Out/2026)}\par\vspace{2pt}
  {\fontsize{7.0}{8.6}\selectfont\setlength{\tabcolsep}{4.5pt}%
\begin{tabularx}{\linewidth}{>{\raggedright\arraybackslash}Xrr}
\toprule
\rowcolor{corprimaria} \textbf{\color{white}Porte Empresarial (Receita Federal / SINAC)} & \textbf{\color{white}CNPJs Ativos no Tocantins} & \textbf{\color{white}Participação no Estoque (\%)} \\
\midrule
MEI & 88.650 & 52,90\% \\
\rowcolor{corfundo} Microempresa (ME) & 56.958 & 33,99\% \\
Pequeno Porte (EPP) & 11.896 & 7,10\% \\
\rowcolor{corfundo} Não Informado / Demais & 10.061 & 6,00\% \\
\bottomrule
\end{tabularx}%
}
  \par\vspace{1.5pt}{\fontsize{6.0}{7.2}\selectfont\color{cinzasub}\textit{Fonte: Receita Federal / SERPRO e SINAC (Optantes Simples Nacional: 164.996 | Tempo Médio de Abertura 2026: 18,5h).}}
\end{minipage}
\par\vspace{0.08cm}

\par\vspace{0.08cm}
\noindent\begin{minipage}[t]{0.482\textwidth}
\par\vspace{0.08cm}
\noindent\begin{minipage}{\linewidth}
  \noindent{\fontsize{7.8}{9.5}\selectfont\bfseries\color{corprimaria}\rule{0.10cm}{0.20cm}\hspace{4pt}Quadro Societário (QSA) por Qualificação (Total: 106.471)}\par\vspace{2pt}
  {\scriptsize\setlength{\tabcolsep}{4.0pt}%
\noindent\resizebox{\linewidth}{!}{%
\begin{tabular}{lrrr}
\toprule
\rowcolor{corprimaria} \textbf{\color{white}Qualificação e Tipo no QSA (RFB)} & \textbf{\color{white}Sócios} & \textbf{\color{white}Empresas} & \textbf{\color{white}Part. \%} \\
\midrule
Sócio-Administrador (Pessoa Física) & 61.828 & 54.177 & 55,33\% \\
\rowcolor{corfundo} Sócio (Pessoa Física) & 18.061 & 13.176 & 16,16\% \\
Diretor (Pessoa Física) & 14.635 & 1.996 & 13,10\% \\
\rowcolor{corfundo} Administrador (Pessoa Física) & 7.706 & 5.304 & 6,90\% \\
Sócio (Pessoa Jurídica) & 4.789 & 3.274 & 4,29\% \\
\rowcolor{corfundo} Conselheiro de Administração (Pessoa Física) & 1.854 & 515 & 1,66\% \\
\bottomrule
\end{tabular}%
}%
}
  \par\vspace{1.5pt}{\fontsize{6.0}{7.2}\selectfont\color{cinzasub}\textit{Fonte: RFB/SERPRO. Gênero QSA: Feminino (32,9\%); Masculino (67,0\%); Não informado (0,0\%).}}
\end{minipage}
\par\vspace{0.08cm}
\end{minipage}\hfill
{\color{cinzaborda}\vrule width 0.4pt}\hfill
\begin{minipage}[t]{0.482\textwidth}
\par\vspace{0.08cm}
\noindent\begin{minipage}{\linewidth}
  \noindent{\fontsize{7.8}{9.5}\selectfont\bfseries\color{corprimaria}\rule{0.10cm}{0.20cm}\hspace{4pt}Perfil Etário do Quadro Societário (QSA — TO)}\par\vspace{2pt}
  {\scriptsize\setlength{\tabcolsep}{4.0pt}%
\noindent\resizebox{\linewidth}{!}{%
\begin{tabular}{lrrr}
\toprule
\rowcolor{corprimaria} \textbf{\color{white}Faixa Etária dos Sócios (QSA)} & \textbf{\color{white}Sócios} & \textbf{\color{white}Empresas} & \textbf{\color{white}Part. \%} \\
\midrule
Até 25 anos & 4.821 & 4.331 & 4,5\% \\
\rowcolor{corfundo} Dos 26 aos 35 anos & 18.579 & 16.230 & 17,4\% \\
Dos 36 aos 49 anos & 42.435 & 30.040 & 39,9\% \\
\rowcolor{corfundo} Dos 50+ & 40.635 & 24.007 & 38,2\% \\
\bottomrule
\end{tabular}%
}%
}
  \par\vspace{1.5pt}{\fontsize{6.0}{7.2}\selectfont\color{cinzasub}\textit{Fonte: Receita Federal do Brasil / SERPRO — Microdados do QSA.}}
\end{minipage}
\par\vspace{0.08cm}
\end{minipage}
\par\vspace{0.08cm}

\bannersecaoperiodico{2}{BALANÇO RETROSPECTIVO (12 MESES) E CENÁRIOS PROSPECTIVOS}{Séries Anualizadas, Consolidação Plurianual, Capilaridade Regional e Perspectivas da Equipe Técnica}
\par\vspace{0.08cm}
\noindent\begin{minipage}{\linewidth}
  \noindent{\fontsize{7.8}{9.5}\selectfont\bfseries\color{corprimaria}\rule{0.10cm}{0.20cm}\hspace{4pt}Balanço Demográfico Empresarial nos Últimos 12 Meses no Tocantins (Aberturas, Baixas e Estoque)}\par\vspace{2pt}
  {\fontsize{7.0}{8.6}\selectfont\setlength{\tabcolsep}{4.5pt}%
\begin{tabularx}{\linewidth}{>{\raggedright\arraybackslash}Xrrrr}
\toprule
\rowcolor{corprimaria} \textbf{\color{white}Competência} & \textbf{\color{white}Novas Inscrições (Aberturas)} & \textbf{\color{white}Baixas / Extinções} & \textbf{\color{white}Saldo Líquido Mensal} & \textbf{\color{white}Estoque de CNPJs Ativos} \\
\midrule
Nov/25 & 2.168 & 1.306 & +862 & 152.819 \\
\rowcolor{corfundo} Dez/25 & 2.007 & 1.441 & +566 & 153.385 \\
Jan/26 & 3.845 & 2.042 & +1.803 & 155.188 \\
\rowcolor{corfundo} Fev/26 & 3.241 & 1.781 & +1.460 & 156.648 \\
Mar/26 & 3.518 & 1.790 & +1.728 & 158.376 \\
\rowcolor{corfundo} Abr/26 & 3.086 & 1.636 & +1.450 & 159.826 \\
Mai/26 & 2.973 & 1.798 & +1.175 & 161.001 \\
\rowcolor{corfundo} Jun/26 & 3.009 & 1.690 & +1.319 & 162.320 \\
Jul/26 & 2.781 & 1.468 & +1.313 & 163.633 \\
\rowcolor{corfundo} Ago/26 & 2.989 & 1.562 & +1.427 & 165.060 \\
Set/26 & 2.617 & 1.416 & +1.201 & 166.261 \\
\rowcolor{corfundo} Out/26 & 2.760 & 1.456 & +1.304 & 167.565 \\
\bottomrule
\end{tabularx}%
}
  \par\vspace{1.5pt}{\fontsize{6.0}{7.2}\selectfont\color{cinzasub}\textit{Fonte: Receita Federal do Brasil / SERPRO.}}
\end{minipage}
\par\vspace{0.08cm}

\vspace{0.03cm}
\begin{multicols}{2}
\secaoperiodico{2.1. Vetores Estruturais e Perspectivas para o Próximo Ciclo}
{\small A consolidação do tempo médio de abertura em 18,5 horas no Mapa de Empresas (REDESIM) e a elevada adesão ao Simples Nacional (164.996 optantes) projetam continuidade do ciclo de formalização empresarial ao longo do próximo exercício.}

\secaoperiodico{2.2. Transmissão Intersetorial e Capilaridade Municipal}
{\small O acompanhamento sob a ótica acumulada em 12 meses evidencia que a dinâmica econômica do Tocantins apoia-se na complementaridade entre os polos regionais consolidados (Palmas, Araguaína, Gurupi, Porto Nacional, Paraíso do Tocantins e Colinas do Tocantins) e os municípios de expansão agroindustrial e logística ao longo do corredor da BR-153 e da Ferrovia Norte-Sul. A depuração de efeitos sazonais de curto prazo permite identificar ganhos persistentes de escala produtiva, formalização empresarial e densidade salarial.}

\secaoperiodico{2.3. Síntese Institucional e Metodológica (Ótica de 12 Meses)}
{\small A leitura conjunta das séries acumuladas em quatro trimestres e das comparações interanuais confirma a solidez dos fundamentos regionais do Tocantins neste eixo temático. O monitoramento contínuo das bases oficiais pelo Jornal Econômico do Tocantins — integrando extração automatizada via APIs públicas, tratamento de microdados e revisão analítica da equipe — permite distinguir oscilações transitórias de mudanças estruturais de tendência, subsidiando decisões de investimento privado e formulação de políticas públicas estaduais e municipais.}
\end{multicols}


\vspace{0.12cm}

% ==============================================================================
% EXPEDIENTE INSTITUCIONAL E NOTA METODOLÓGICA
% ==============================================================================
\noindent\begin{tikzpicture}
  \node[
    fill=corfundo,
    draw=cinzaborda,
    line width=0.5pt,
    rounded corners=2pt,
    inner sep=6pt,
    text width=\textwidth-13pt,
    font=\scriptsize\color{cinzasub}
  ] {
    \textbf{\color{corprimaria}Expediente e Nota Metodológica:} Integra os registros de CNPJs ativos por porte (MEI, ME, EPP e Demais), optantes pelo Simples Nacional (SINAC), tempos de abertura (Mapa de Empresas) e o Quadro de Sócios e Administradores (QSA — qualificação, faixa etária e sexo).\par\vspace{2pt}
    \textbf{\color{corprimaria}Equipe Responsável:} Lorenzo Costa | Graduado em Economia — Jornal Econômico do Tocantins \enspace$\vert$\enspace \textbf{\color{corprimaria}Publicação Institucional:} Núcleo de Conjuntura Econômica — \textit{Jornal Econômico do Tocantins}.
  };
\end{tikzpicture}

\end{document}
